\section{Additional Mini-Projects}
\label{extra_mini_projects}

If you feel like challenging yourself, getting some additional practice, or just feel like tinkering a little more, here are some additional circuit ideas you could pursue.

\bigskip
\textbf{Mini-Project 3: An Analog-to-Digital Converter}

In Lab~\ref{lab_op-amps}, you built a D-to-A converter.  Here, you will build the opposite: an A-to-D converter, which takes a single, variable voltage input and converts it to a series of discrete digital voltages representing a base-two number.  This would be the kind of circuit you would need to turn a signal from a microphone into a digital sound file, or to turn an oscilloscope's input signal into a value to graph on the display and analyze.

You can choose whether to make a two-bit converter for 0 to 3 volts, or a three-bit converter for 0 to 7 volts.  
For the two-bit version, your output will be a set of two digital signals representing the numbers 00, 01, 10, and 11.  (0.00--0.99 volts gives digital 00; 1.00--1.99 volts gives digital 01, etc.)  Here's how you'll go about it: start by comparing the input signal to a 2-volt reference and set the twos bit accordingly.  If $V_{\rm in}<2$~volts, then you'll compare the input to a 1-volt reference to set the ones bit.  
But if $V_{\rm in}>2$, then you'll need to subtract off 2~volts from $V_{\rm in}$ first, and \textit{then} set the ones bit by comparing $V_{\rm in}-2$~volts to the 1-volt reference.  Build and test your circuit and draw a circuit diagram for it in your notebook.

\begin{itemize}
\item You can see how the scheme used here would become pretty complicated for an 8-bit or 16-bit converter.  Think about how you could design a 4-bit converter to use as few components as possible.  In particular, there are surely smart ways to take advantage of the comparator's flexible open-collector output.

\item Feel free to use a different comparator than the LM311; the LM393 has the advantage of two comparators on a single chip. 

\item This multi-stage compare-and-subtract strategy is only \textit{one} way to build an A-to-D converter, and not necessarily the best way.  You can read more about it in Horowitz and Hill's \textit{Art of Electronics}.
\end{itemize}

\bigskip
\textbf{Mini-Project 4: A Constant Current Supply}

Although most power supplies work by supplying a constant voltage, one occasionally needs a power supply to provide a constant current.  Some examples would be:
\begin{itemize}
\item To account for line voltage that varies over time, for instance to provide a fixed current to a bank of LED lights to maintain constant brightness and avoid damaging them by an overvoltage.
\item To account for a varying load resistance, for instance to provide a fixed current to a motor on a model train engine, even when the train is far from the power supply, so that the length of track adds additional resistance to the circuit.
\item To  guard against an overcurrent in the load due to a short circuit. The big DC power supplies that we used in Lab~\ref{lab_power_dissipation} did this, allowing us to adjust the current knobs to set an effective current limit for the voltage supply.
\end{itemize}

Build a circuit that provides a fixed current of 10~mA~DC to a load that can vary from $200\ohm$ to $1\kohm$, starting with a voltage source that can vary from $+12$~VDC to $+20$~VDC.  Your basic strategy will be to use an op-amp to supply the output voltage, but to monitor the output current, using the current as the basis for negative feedback to the op-amp.

Modify your circuit so that you can adjust the output current to anywhere from 0 to 20~mA by turning a potentiometer---to make that model train go faster or slower.

If you know how to use bipolar junction transistors (Lab~\ref{lab_bjt}), you can build a higher-current version of your circuit that delivers up to 1~amp to a sufficiently low-impedance load. 
 


